\begin{abstract}
Coding agents run inside harnesses that execute their tool calls, and those tools fail: commands error, reads return nothing or stale data, and calls that succeeded report failure.
We present Chaos Monkey for Coding Agents (CMCA), a fault-injection benchmark that applies one seeded fault model inside unmodified production harnesses at two layers.
On the service track, an MCP proxy in front of mock services delivers identical bytes to every harness.
On the native track, injectors loaded through each harness's own hook or extension interface fault its built-in read, edit, search and shell tools, and comparisons are restricted to fault types that a conformance suite shows each harness delivers.
A fixed-model axis runs every harness with the same model, so differences are attributable to the harness.
Besides task success under faults, CMCA measures whether the agent reports its outcome honestly, whether stale data reaches its output, and what recovery costs.
\ifcmcaresults
On \cmcaTrackSTasks{} service and \cmcaTrackNFixedTasks{} native tasks with one model in every harness, the harness that completes the most faulted tasks is \cmcaTrackSLeader{} on the service track (\cmcaTrackSLeaderSucceeded{} of \cmcaTrackSLeaderSent{}) and \cmcaTrackNFixedLeader{} on the native track (\cmcaTrackNFixedLeaderSucceeded{} of \cmcaTrackNFixedLeaderSent{}), but no paired difference survives multiple-comparison correction, and among failed faulted trials agents report success in \cmcaTrackSOverclaimMin{} to \cmcaTrackNFixedOverclaimMax{} of cases.
\fi
Code, tasks, trajectories and fault logs are public.
\end{abstract}
